% !TeX encoding = UTF-8
% Formatting-only conversion; see reports for source provenance.
\SourceChapter{1}{فصل اول: کلیات پژوهش}
% SOURCE Introduction_Chapter1.txt:1-1 [chapter-title-in-wrapper]


% SOURCE Introduction_Chapter1.txt:2-2 [spacing]


% SOURCE Introduction_Chapter1.txt:3-3 [heading]
\SourceHeading{1}{۱{-}۱. مقدمه}

% SOURCE Introduction_Chapter1.txt:4-4 [paragraph]
با توسعه روزافزون فناوری‌های سلامت الکترونیک \lr{(E{-}Health)}، پزشکی از راه دور \lr{(Telemedicine)} و اینترنت اشیای پزشکی \lr{(IoMT)}، تبادل الکترونیکی تصاویر تشخیصی نظیر تصاویر تصویربرداری تشدید مغناطیسی \lr{(MRI)}، توموگرافی کامپیوتری \lr{(CT)}، رادیوگرافی قفسه سینه \lr{(X{-}Ray)} و تصویربرداری شبکیه چشم به بخشی لاینفک از فرآیندهای نوین درمانی تبدیل شده است \lr{[1, 2]}. این داده‌های تصویری نه تنها اطلاعات آناتومیک بسیار حساسی را جهت تشخیص بیماری‌ها ارائه می‌دهند، بلکه حاوی فراداده‌های هویتی و پرونده‌های بالینی بیماران نیز هستند \lr{[7, 9]}. بنابراین، حفظ محرمانگی، اصالت و یکپارچگی این داده‌ها هنگام انتقال بر روی کانال‌های ارتباطی عمومی و ناامن شهری و بین‌المللی، یکی از چالش‌های بنیادین در فضای سلامت دیجیتال محسوب می‌گردد \lr{[2, 10]}.\par

% SOURCE Introduction_Chapter1.txt:5-5 [spacing]


% SOURCE Introduction_Chapter1.txt:6-6 [paragraph]
روش‌های کلاسیک و متعارف رمزنگاری نظیر استاندارد رمزنگاری پیشرفته \lr{(AES)} و \lr{RSA}، اگرچه برای داده‌های متنی عملکرد استواری دارند، اما در مواجهه با تصاویر پزشکی با محدودیت‌های شدیدی روبه‌رو هستند. تصاویر پزشکی دارای ویژگی‌های منحصربه‌فردی همچون حجم سنگین داده، افزونگی اطلاعاتی بسیار بالا و همبستگی شدید میان پیکسل‌های همجوار در جهت‌های افقی، عمودی و قطری می‌باشند \lr{[1, 3, 6]}. الگوریتم‌های سنتی متنی به دلیل پیچیدگی محاسباتی بالا و نرخ پردازش پایین، توانایی رمزنگاری زمان‌واقعی \lr{(Real{-}time)} این دادگان حجیم را نداشته و قادر به شکستن کامل همبستگی پیکسل‌های مجاور در تصاویر دادگان بالینی نیستند \lr{[3, 6]}.\par

% SOURCE Introduction_Chapter1.txt:7-7 [spacing]


% SOURCE Introduction_Chapter1.txt:8-8 [paragraph]
برای غلبه بر این چالش‌ها، استفاده از سیستم‌های غیرخطی دینامیکی و نظریه آشوب به عنوان یکی از کارآمدترین راهکارهای رمزنگاری داده‌های چندرسانه‌ای مطرح گردیده است. سیستم‌های آشوبی به دلیل ویژگی‌هایی نظیر حساسیت شدید به شرایط اولیه و پارامترهای کنترل، عدم دوره‌ای بودن \lr{(Non{-}periodicity)}، شبه‌تصادفی بودن فوق‌العاده و شبه‌نویز بودن خروجی، بستری ایده‌آل برای دو فاز اصلی رمزنگاری یعنی آشفته‌سازی \lr{(Confusion)} و انتشار \lr{(Diffusion)} فراهم می‌سازند \lr{[1, 3, 6]}. با این حال، سیستم‌های آشوبی کم‌بعد (مانند نگاشت لوجستیک یک‌بعدی) به دلیل فضای کلید محدود و امکان بازسازی فاز، در برابر حملات پیشرفته سایبری آسیب‌پذیرند \lr{[3]}. از این رو، به‌کارگیری سیستم‌های ابرآشوبی چندبعدی \lr{(Hyperchaotic Systems)} با بیش از یک توان نمای لیاپانوف مثبت \lr{(Positive Lyapunov Exponent)}، پیچیدگی و امنیت رمزنگاری را به طرز چشمگیری ارتقا می‌دهد \lr{[1, 10]}. سیستم ابرآشوبی ۵ بعدی پیشنهادی در این پژوهش، به دلیل دارا بودن توان‌های لیاپانوف مثبت متعدد و رفتارهای نوسانی پیچیده، فضای کلیدی فراتر از \lr{2}\SourceGlyph{005E}\lr{500} ایجاد نموده و امکان هرگونه بازسازی فضای فاز \lr{(Phase Space Reconstruction)} یا حملات جستجوی فراگیر \lr{(Brute{-}force)} را سلب می‌کند، ضمن آنکه بستر مناسبی را برای دستیابی به آنتروپی آرمانی (\lr{7.9985} بیت) و شاخص‌های تفاضلی استاندارد (\lr{NPCR \SourceGlyph{003E} 99.62\%} و \lr{UACI} \SourceGlyph{007E} \lr{33.51\%}) فراهم می‌نماید \lr{[1, 3, 6]}.\par

% SOURCE Introduction_Chapter1.txt:9-9 [spacing]


% SOURCE Introduction_Chapter1.txt:10-10 [paragraph]
از سوی دیگر، یکی از نقاط ضعف اصلی الگوریتم‌های آشوبی سنتی، عدم وابستگی مستقیم کلید رمزنگاری به محتوای تصویر ورودی است. در صورتی که کلید رمزنگاری مستقل از تصویر باشد، سیستم در برابر حملات متن‌واضح برگزیده \lr{(Chosen{-}Plaintext Attack)} و متن‌واضح معلوم \lr{(Known{-}Plaintext Attack)} آسیب‌پذیر خواهد بود \lr{[1, 5]}. برای رفع این خلاء، در این پژوهش از شبکه‌های عصبی عمیق به‌ویژه معماری \lr{U{-}Net} جهت قطعه‌بندی معنایی نواحی حساس پزشکی \lr{(ROI)} و استخراج خصوصیات آماری وابسته به محتوا استفاده می‌شود \lr{[1, 2]}. انتخاب معماری \lr{U{-}Net} به دلیل دارا بودن ساختار تقارنی رمزگذار{-}رمزگشا \lr{(Encoder{-}Decoder)} و اتصالات پرشی \lr{(Skip Connections)}، امکان استخراج دقیق مرزهای آناتومیک حساس (نظیر عروق شبکیه و تومورهای مغزی) را حتی با حجم داده‌های آموزشی محدود فراهم می‌سازد و با بار محاسباتی بهینه، تولید کلیدهای پویای وابسته به تصویر را جهت ارتقای امنیت سیستم تضمین می‌نماید \lr{[1, 2]}.\par

% SOURCE Introduction_Chapter1.txt:11-11 [spacing]


% SOURCE Introduction_Chapter1.txt:12-12 [paragraph]
علاوه بر رمزنگاری، صیانت از فراداده‌های حساس بیمار (مانند مشخصات هویتی و سوابق بالینی) بدون ایجاد اعوجاج در کیفیت تشخیصی تصویر، یکی دیگر از ارکان این چارچوب است \lr{[4, 7, 8]}. در این سامانه، فراداده‌های متنی ابتدا با استفاده از الگوریتم \lr{zlib} فشرده‌سازی شده تا حجم پیام به حداقل برسد، سپس با بهره‌گیری از الگوریتم پنهان‌نگاری آگاه به برجستگی بصری \lr{(Saliency{-}aware Steganography)}، فراداده فشرده‌شده صرفاً در نواحی کم‌اهمیت دیداری و پس‌زمینه غیرحساس تصویر پزشکی جای‌دهی می‌شود \lr{[4, 5, 7]}. این رویه هوشمندانه باعث می‌شود که نواحی تشخیصی اصلی \lr{(ROI)} کاملاً دست‌نخورده باقی مانده و کیفیت دیداری تصویر با معیار \lr{PSNR} بالاتر از \lr{53} دسی‌بل و \lr{SSIM} نزدیک به \lr{0.998} حفظ گردد و سیستم در برابر حملات آماری و دیداری کاملاً مقاوم باشد \lr{[7, 8]}.\par

% SOURCE Introduction_Chapter1.txt:13-13 [spacing]


% SOURCE Introduction_Chapter1.txt:14-14 [paragraph]
تلفیق یادگیری عمیق، سیستم‌های ابرآشوبی ۵ بعدی و تکنیک‌های پنهان‌نگاری فراداده مبتنی بر برجستگی بصری \lr{(Saliency)}، چارچوبی جامع، نوین و زمان‌واقعی (با زمان اجرای متوسط \lr{2.93} ثانیه) برای حفاظت همزمان از محرمانگی تصویر و فراداده‌های بیمار ارائه می‌دهد \lr{[1, 4, 7]}.\par

% SOURCE Introduction_Chapter1.txt:15-15 [spacing]


% SOURCE Introduction_Chapter1.txt:16-16 [paragraph]
در ادامه این فصل، ابتدا به بیان دقیق مسئله و اهداف اصلی و فرعی پژوهش پرداخته شده، سپس پرسش‌ها، فرضیات و جنبه‌های نوآوری سیستم پیشنهادی تبیین می‌گردند. در نهایت با مرور اهمیت، حدود و قلمرو پژوهش و ارائه تعاریف نظری و عملیاتی متغیرها، نقشه راه کل پایان‌نامه ترسیم خواهد شد.\par
% SOURCE Problem_Statement_Chapter1.txt:1-1 [heading]
\SourceHeading{1}{۱{-}۲. بیان مسئله}

% SOURCE Problem_Statement_Chapter1.txt:2-2 [spacing]


% SOURCE Problem_Statement_Chapter1.txt:3-3 [paragraph]
با گسترش روزافزون سامانه‌های سلامت الکترونیک \lr{(E{-}Health)}، بیمارستان‌های هوشمند، خدمات ابری پزشکی و بسترهای سامانه بایگانی و تبادل تصاویر \lr{(PACS)}، تبادل الکترونیکی دادگان تشخیصی مبتنی بر تصویر نظیر تصاویر تصویربرداری تشدید مغناطیسی \lr{(MRI)}، توموگرافی کامپیوتری \lr{(CT)}، رادیوگرافی قفسه سینه \lr{(X{-}Ray)} و تصویربرداری شبکیه چشم به امری اجتناب‌ناپذیر و زیربنایی تبدیل شده است \lr{[1, 2]}. این داده‌های تصویری نه تنها اطلاعات آناتومیک بسیار حساسی را جهت تشخیص دقیق بیماری‌ها، تومورها و عروق ارائه می‌دهند، بلکه حاوی فراداده‌های هویتی، پرونده‌های بالینی و شماره‌های شناسایی بیماران نیز می‌باشند \lr{[7, 9]}. بر اساس گزارش‌های معتبر بین‌المللی در حوزه امنیت اطلاعات سلامت، نشت داده‌های پزشکی و حملات سایبری به بسترهای بیمارستانی در سال‌های اخیر رشد فزاینده‌ای داشته است؛ به‌گونه‌ای که دادگان پزشکی به یکی از ارزشمندترین اهداف در بازار دارک‌وب تبدیل شده‌اند \lr{[2, 10]}. این امر صیانت از یکپارچگی، محرمانگی و حریم خصوصی داده‌های بالینی را بر اساس قوانین حاکمیتی بین‌المللی نظیر \lr{HIPAA} (قانون انتقال و پاسخگویی بیمه سلامت) و \lr{GDPR} (مقررات عمومی حفاظت از داده اتحادیه اروپا) به یک الزام حقوقی، اخلاقی و فنی مبرم تبدیل می‌سازد \lr{[2, 9]}.\par

% SOURCE Problem_Statement_Chapter1.txt:4-4 [spacing]


% SOURCE Problem_Statement_Chapter1.txt:5-5 [paragraph]
انتقال امن تصاویر پزشکی در بسترهای سلامت الکترونیک و اینترنت اشیای پزشکی \lr{(IoMT)} با سه چالش اساسی و ساختاری روبه‌رو است که راهکارهای تک‌بعدی و سنتی موجود قادر به حل همزمان آن‌ها نیستند \lr{[1, 2, 8]}:\par

% SOURCE Problem_Statement_Chapter1.txt:6-6 [spacing]


% SOURCE Problem_Statement_Chapter1.txt:7-7 [paragraph]
۱. چالش همبستگی پیکسلی و حجم بالا: تصاویر تشخیصی پزشکی دارای مقادیر تشابه بسیار بالایی میان پیکسل‌های مجاور هستند. الگوریتم‌های متعارف رمزنگاری الگوی بصری تصویر را به خوبی پنهان نمی‌سازند یا سرعت بسیار پایینی در پردازش این دادگان حجیم دارند \lr{[3, 6]}. الگوریتم‌های ایستا و متنی سنتی نظیر \lr{AES} (در حالت‌های عملیاتی متعارف) و \lr{RSA}، به دلیل فرض ساختار متنی یک‌بعدی، قادر به شکستن همبستگی شدید پیکسلی (که در سه جهت افقی، عمودی و قطری مقادیری نزدیک به ۱ دارد) نبوده و به دلیل پیچیدگی محاسباتی بالا، در مواجهه با ماتریس‌های حجیم تصاویر پزشکی تشخیصی باعث ایجاد گلگاه \lr{(Bottleneck)} پردازشی در تبادلات زمان‌واقعی \lr{(Real{-}time)} می‌گردند \lr{[1, 3, 6]}.\par

% SOURCE Problem_Statement_Chapter1.txt:8-8 [spacing]


% SOURCE Problem_Statement_Chapter1.txt:9-9 [paragraph]
۲. آسیب‌پذیری در برابر حملات متن‌واضح (ایستا بودن کلید): در بسیار از سامانه‌های آشوبی موجود، کلید رمزنگاری برای تصاویر مختلف یکسان یا مستقل از محتوای تصویر است. این امر به نفوذگران اجازه می‌دهد با ارسال تصاویر آزمایشی، ساختار کلید را بازسازی و امنیت سیستم را خنثی کنند \lr{[1, 5]}. علاوه بر این، استفاده از سیستم‌های آشوبی کم‌بعد (مانند نگاشت‌های لوجستیک یا خیمه‌ای یک‌بعدی و دوبعدی) به دلیل محدودیت فضای کلید (کمتر از \lr{2}\SourceGlyph{005E}\lr{128})، دوره تناوب کوتاه توالی‌ها و امکان بازسازی فضای فاز \lr{(Phase Space Reconstruction)} توسط الگوریتم‌های پنهان‌کاوی دینامیکی، سیستم را در برابر حملات متن‌واضح برگزیده \lr{(CPA)} و متن‌واضح معلوم \lr{(KPA)} به شدت آسیب‌پذیر می‌سازد \lr{[1, 3, 6]}. در چنین شرایطی، کوچک‌ترین تغییر در تصویر ورودی منجر به تغییر محسوس در کلید نشده و عدم وجود اثر بهمنی \lr{(Avalanche Effect)} کامل، امنیت سامانه را مختل می‌کند \lr{[1, 5]}.\par

% SOURCE Problem_Statement_Chapter1.txt:10-10 [spacing]


% SOURCE Problem_Statement_Chapter1.txt:11-11 [paragraph]
۳. چالش افشای فراداده‌های حساس بیمار و تخریب کیفیت تشخیصی: فراداده‌های هویتی، پرونده پزشکی و نظرات تشخیصی پزشک معمولاً به صورت جداگانه ارسال شده یا هنگام پنهان‌نگاری در تصویر، باعث ایجاد اعوجاج دیداری و کاهش کیفیت بالینی تصویر پزشکی می‌شوند \lr{[4, 7, 8]}. روش‌های کلاسیک پنهان‌نگاری متنی و تصویری (مانند جایگزینی کم‌ارزش‌ترین بیت \lr{LSB}) با اعمال تغییرات غیرتطبیقی در تمامی سطوح تصویر، علاوه بر ظرفیت محدود در گنجاندن فراداده‌های حجیم، تغییرات آماری قابل کشفی در هیستوگرام تصویر ایجاد کرده و توسط ابزارهای پنهان‌کاوی \lr{(Steganalysis)} کشف می‌گردند؛ مهم‌تر آن‌که جای‌دهی ناآگاهانه داده در نواحی حساس تشخیصی \lr{(ROI)}، آسیب‌های جبران‌ناپذیری به دقت تشخیص‌پذیری ضایعات و آناتومی بالینی وارد می‌آورد \lr{[4, 7]}.\par

% SOURCE Problem_Statement_Chapter1.txt:12-12 [spacing]


% SOURCE Problem_Statement_Chapter1.txt:13-13 [paragraph]
مسئله اصلی این پژوهش، چگونگی طراحی و پیاده‌سازی یک چارچوب امنیتی هیبریدی و یکپارچه است که بتواند با بهره‌گیری از شبکه عصبی عمیق \lr{U{-}Net}، ویژگی‌های آماری ناحیه حساس تصویر پزشکی \lr{(ROI)} را به صورت کلید پویای وابسته به تصویر استخراج کرده و سپس با اعمال یک سیستم ابرآشوبی ۵ بعدی و الگوریتم‌های پیشرفته انتشار \lr{DNA}، تصویر را با آنتروپی آرمانی رمزنگاری نماید \lr{[1, 2]}. در این فرایند، استخراج خصوصیات آماری و بافتی \lr{ROI} توسط شبکه \lr{U{-}Net} و ترکیب آن با تابع هش \lr{SHA{-}256}، کلیدی کاملاً منحصر‌به‌فرد تولید می‌کند به‌گونه‌ای که با کوچک‌ترین تغییر حتی در یک پیکسل از تصویر ورودی، خروجی کلید به صورت بنیادی تغییر یافته و اثر بهمنی کاملی ایجاد می‌گردد تا امکان هرگونه حمله \lr{CPA} و \lr{KPA} سلب شود \lr{[1, 2, 5]}. همزمان، سیستم ابرآشوبی ۵ بعدی به دلیل دارا بودن بیش از یک توان نمای لیاپانوف مثبت، فضای کلیدی فراتر از \lr{2}\SourceGlyph{005E}\lr{500} و پیچیدگی دینامیکی فوق‌العاده‌ای ایجاد کرده و با اجرای آشفته‌سازی زیگ‌زاگی \lr{(Zig{-}zag Scrambling)} و فازهای پویای انتشار \lr{DNA}، همبستگی پیکسلی را به صفر نزدیک ساخته و آنتروپی اطلاعات را به حد آرمانی \lr{7.9985} بیت می‌رساند \lr{[1, 3, 6]}.\par

% SOURCE Problem_Statement_Chapter1.txt:14-14 [spacing]


% SOURCE Problem_Statement_Chapter1.txt:15-15 [paragraph]
همزمان، این سامانه باید فراداده‌های حساس بیمار را با استفاده از فشرده‌سازی \lr{zlib} و تکنیک پنهان‌نگاری مبتنی بر برجستگی بصری \lr{(Saliency Detection)}، در نواحی کم‌اهمیت تصویر جای‌دهی کند به‌گونه‌ای که هیچ‌گونه افت کیفیت تشخیصی \lr{(PSNR \SourceGlyph{003E} 50 dB)} رخ ندهد \lr{[4, 7]}. با فشرده‌سازی فراداده توسط \lr{zlib} و شناساگر برجستگی بصری، اطلاعات حساس هویتی بیمار صرفاً در نواحی کم‌توجه دیداری و پس‌زمینه تصویر (خارج از محدوده \lr{ROI}) جای‌دهی می‌شود؛ بدین ترتیب، ظرفیت گنجاندن داده افزایش یافته، امنیت فراداده در برابر ابزارهای پنهان‌کاوی تضمین شده و کیفیت دیداری و آناتومیک تصویر تشخیصی با \lr{PSNR} بالای \lr{50} دسی‌بل و \lr{SSIM} نزدیک به \lr{1} به صورت کامل دست‌نخورده باقی می‌ماند \lr{[4, 7, 8]}.\par

% SOURCE Problem_Statement_Chapter1.txt:16-16 [spacing]


% SOURCE Problem_Statement_Chapter1.txt:17-17 [paragraph]
سؤال محوری این پژوهش بدین شرح فرمول‌بندی می‌شود: چگونه می‌توان قطعه‌بندی عمیق با \lr{U{-}Net}، رمزنگاری ابرآشوبی ۵ بعدی همراه با انتشار پویای \lr{DNA}، و پنهان‌نگاری آگاه به برجستگی بصری را در یک زنجیره پردازشی واحد گرد آورد تا ضمن دستیابی به بالاترین سطح محرمانگی، آنتروپی آرمانی و مقاومت کامل در برابر حملات سایبری، کیفیت تشخیصی بالینی تصویر و سرعت اجرای پردازش زمان‌واقعی برای سامانه‌های سلامت الکترونیک و \lr{IoMT} تضمین گردد؟\par

% SOURCE Problem_Statement_Chapter1.txt:18-18 [spacing]

% SOURCE Objectives_Chapter1.txt:1-1 [heading]
\SourceHeading{1}{۱{-}۳. اهداف پژوهش}

% SOURCE Objectives_Chapter1.txt:2-2 [spacing]


% SOURCE Objectives_Chapter1.txt:3-3 [paragraph]
اهداف این پژوهش در دو سطح «هدف اصلی (کلان)» و «اهداف فرعی و کمّی» سازمان‌دهی شده‌اند. این اهداف به گونه‌ای تدوین گردیده‌اند که تناظر یک‌به‌یک و مستقیم با چالش‌های سه‌گانه تبیین‌شده در بخش بیان مسئله (شامل چالش کلید ایستا، همبستگی شدید پیکسلی و افشای فراداده‌های حساس بیمار) داشته و سنجه‌های ارزیابی آنها بر اساس معیارهای سنجش‌پذیر علمی (معیارهای \lr{SMART}) و استانداردهای رایج در امن‌سازی دادگان زیست‌پزشکی اعتبارسنجی شوند \lr{[1, 2, 7]}.\par

% SOURCE Objectives_Chapter1.txt:4-4 [spacing]


% SOURCE Objectives_Chapter1.txt:5-5 [heading]
\SourceHeading{2}{۱{-}۳{-}۱. هدف اصلی}

% SOURCE Objectives_Chapter1.txt:6-6 [paragraph]
طراحی، پیاده‌سازی و ارزیابی جامع یک چارچوب امنیتی یکپارچه و چندلایه‌ای جهت رمزنگاری مقاوم و پنهان‌نگاری تطبیقی تصاویر پزشکی بر پایه یادگیری عمیق \lr{(U{-}Net)}، سیستم‌های ابرآشوبی ۵ بعدی و نگاشت‌های پویای \lr{DNA}، به منظور ارتقای محرمانگی و صیانت از حریم خصوصی بیماران در سامانه‌های سلامت همراه و اینترنت اشیای پزشکی \lr{[1, 2, 7]}.\par

% SOURCE Objectives_Chapter1.txt:7-7 [spacing]


% SOURCE Objectives_Chapter1.txt:8-8 [paragraph]
در چارچوب این هدف کلان، سیستم پیشنهادی می‌کوشد تا بدون دستکاری یا ایجاد اعوجاج دیداری در ناحیه حساس تشخیصی \lr{(ROI)}، سه لایه امنیتی مکمل را متصل سازد: لایه اول، استخراج ویژگی و کلیدسازی پویا بر پایه هوش مصنوعی؛ لایه دوم، آشفته‌سازی زیگ‌زاگی و انتشار ابرآشوبی مبتنی بر سیستم ۵ بعدی و قوانین پویای \lr{DNA}؛ و لایه سوم، پنهان‌نگاری فراداده‌های بیمار مبتنی بر برجستگی بصری \lr{(Saliency Detection)}. خروجی نهایی این هدف، دست‌یابی به سامانه‌ای جامع، پرسرعت و سبک با قابلیت پیاده‌سازی روی تجهیزات لبه \lr{(Edge Computing)} در بسترهای پزشکی از راه دور \lr{(Telemedicine)} و بیمارستان‌های هوشمند است \lr{[1, 2, 8]}.\par

% SOURCE Objectives_Chapter1.txt:9-9 [spacing]


% SOURCE Objectives_Chapter1.txt:10-10 [heading]
\SourceHeading{2}{۱{-}۳{-}۲. اهداف فرعی و کمّی}

% SOURCE Objectives_Chapter1.txt:11-11 [spacing]


% SOURCE Objectives_Chapter1.txt:12-12 [paragraph]
۱. استخراج دقیق ناحیه حساس تشخیصی \lr{(ROI)} از تصاویر پزشکی با استفاده از شبکه \lr{U{-}Net} و محاسبه تابع هش \lr{SHA{-}256} بر روی ویژگی‌های آماری آن جهت تولید کلیدهای رمزنگاری کاملاً پویا و وابسته به محتوا \lr{[1, 2]}.\par

% SOURCE Objectives_Chapter1.txt:13-13 [paragraph]
توضیح و بسط تفصیلی: این هدف فرعی مستقیماً در پاسخ به چالش کلیدهای ایستا و آسیب‌پذیری سیستم در برابر حملات متن‌واضح برگزیده \lr{(CPA)} و متن‌واضح معلوم \lr{(KPA)} تعریف شده است \lr{[1, 5]}. شبکه عصبی عمیق \lr{U{-}Net} با بهره‌گیری از معماری متقارن انقباضی{-}انبساطی و اتصالات پرشی \lr{(Skip Connections)}، نواحی حساس آناتومیک (مانند عروق شبکیه در دادگان \lr{DRIVE/RITE} یا بافت‌های توموری در دادگان \lr{BraTS}) را با دقت بالا قطعه‌بندی می‌کند \lr{[1, 2]}. سپس با استخراج ویژگی‌های آماری ناحیه \lr{ROI} (شامل میانگین، واریانس، گشتاورها و مقادیر پیکسل‌ها) و اعمال تابع هش ایمن \lr{SHA{-}256}، کلید اولیه‌ای به طول ۲۵۶ بیت تولید می‌شود. این کلید برای تنظیم شرایط اولیه و پارامترهای کنترل سیستم ۵ بعدی ابرآشوبی به کار می‌رود؛ در نتیجه، حتی با کوچک‌ترین تغییر در یک پیکسل از تصویر ورودی، کلید رمزنگاری به صورت بنیادی تغییر یافته و اثر بهمنی \lr{(Avalanche Effect)} کامل برقرار می‌گردد \lr{[1, 2]}.\par

% SOURCE Objectives_Chapter1.txt:14-14 [spacing]


% SOURCE Objectives_Chapter1.txt:15-15 [paragraph]
۲. توسعه و تحلیل رفتار دینامیکی یک سیستم ابرآشوبی ۵ بعدی جدید با توان‌های لیاپانوف مثبت متعدد جهت ایجاد توالی‌های شبه‌تصادفی با امنیت فوق‌العاده بالا \lr{[1, 6]}.\par

% SOURCE Objectives_Chapter1.txt:16-16 [paragraph]
توضیح و بسط تفصیلی: این هدف فرعی به منظور برطرف ساختن محدودیت‌های سیستم‌های آشوبی کم‌بعد (۱\lr{D/2D}) تدوین گردیده است \lr{[3, 6]}. سیستم‌های کم‌بعد به دلیل فضای کلید کوچک (کمتر از \lr{2}\SourceGlyph{005E}\lr{128}) و دوره تناوب کوتاه، در برابر بازسازی فضای فاز و حملات جستجوی فراگیر \lr{(Brute{-}force)} آسیب‌پذیر هستند \lr{[3]}. در این پژوهش، با ارتقای ابعاد سیستم به ۵ بعد و فرمول‌بندی معادلات دیفرانسیل غیرخطی کوپل‌شده، سیستمی با بیش از یک توان نمای لیاپانوف مثبت \lr{(Multiple Positive Lyapunov Exponents)} توسعه داده می‌شود \lr{[1]}. این سیستم رفتاری بسیار پیچیده و غیرقابل‌پیش‌بینی ارائه داده و فضای کلیدی فراتر از \lr{2}\SourceGlyph{005E}\lr{500} ایجاد می‌نماید که هرگونه حمله بازسازی الگوریتمی و تحلیل فاز را غیرممکن می‌سازد \lr{[1, 6]}.\par

% SOURCE Objectives_Chapter1.txt:17-17 [spacing]


% SOURCE Objectives_Chapter1.txt:18-18 [paragraph]
۳. اجرای فازهای آشفته‌سازی زیگ‌زاگی \lr{(Zig{-}zag Scrambling)} و انتشار پویای \lr{DNA} \lr{(Dynamic DNA Flip/XOR)} برای شکستن کامل همبستگی پیکسل‌ها و دستیابی به آنتروپی اطلاعات بالای \lr{7.998} بیت \lr{[1]}.\par

% SOURCE Objectives_Chapter1.txt:19-19 [paragraph]
توضیح و بسط تفصیلی: این هدف فرعی به چالش همبستگی پیکسلی شدید (مقادیر همبستگی افقی، عمودی و قطری نزدیک به ۱) در تصاویر تشخیصی پزشکی پاسخ می‌دهد \lr{[1, 3]}. در مرحله آشفته‌سازی، با الگوریتم پیمایش زیگ‌زاگی مبتنی بر توالی‌های شبه‌تصادفی سیستم ۵ بعدی، موقعیت مکانی پیکسل‌ها کاملاً برهم زده می‌شود تا همبستگی مکانی خنثی گردد \lr{[1]}. در مرحله انتشار، پیکسل‌ها به توالی‌های نوکلئوتیدی \lr{DNA} (قواعد رمزنگاری \lr{A, C, G, T}) تبدیل شده و تحت عملیات چرخش پویای \lr{DNA} \lr{(Dynamic DNA Flip)} و تفاضل/\lr{XOR} آشوبی قرار می‌گیرند \lr{[1]}. خروجی این دو فاز متوالی، دستیابی به آنتروپی شانون بسیار بالا (مقدار عددی \lr{7.9985} بیت در برابر مقدار آرمانی ۸ بیت) و همبستگی پیکسلی نزدیک به صفر در تصویر رمزشده است که تصویر را به یک نویز کامل تبدیل می‌کند \lr{[1, 3]}.\par

% SOURCE Objectives_Chapter1.txt:20-20 [spacing]


% SOURCE Objectives_Chapter1.txt:21-21 [paragraph]
۴. طراحی ماژول پنهان‌نگاری فراداده بیمار در نواحی کم‌توجه بصری \lr{(Saliency Detection)} با فشرده‌سازی \lr{zlib} جهت دستیابی به کیفیت دیداری فوق‌العاده با \lr{PSNR} بالاتر از \lr{50} دسی‌بل و \lr{SSIM} نزدیک به \lr{1} \lr{[4, 7]}.\par

% SOURCE Objectives_Chapter1.txt:22-22 [paragraph]
توضیح و بسط تفصیلی: این هدف فرعی پاسخگوی چالش سوم، یعنی افشای فراداده‌های حساس بیمار و جلوگیری از تخریب کیفیت تشخیصی تصویر بالینی است \lr{[4, 7, 8]}. فراداده‌های متنی و هویتی بیمار (نظیر کد ملی، پرونده پزشکی و نظر تشخیصی) ابتدا با الگوریتم \lr{zlib} فشرده‌سازی بی‌زیان شده تا حجم داده کاهش یابد \lr{[4]}. سپس با استفاده از نقشه برجستگی بصری \lr{(Saliency Map)}، نواحی بی‌اهمیت تشخیصی (مانند پس‌زمینه تیره در تصاویر \lr{MRI} و \lr{CT}) شناسایی می‌شوند. فراداده فشرده‌شده صرفاً در این نواحی کم‌توجه جای‌دهی می‌گردد؛ به گونه‌ای که ناحیه اصلی \lr{ROI} کاملاً دست‌نخورده باقی مانده و شاخص نسبت سیگنال به نویز بیشینه \lr{(PSNR)} به بالای ۵۰ دسی‌بل و شاخص شباهت ساختاری \lr{(SSIM)} به بالای \lr{0.998} دست می‌یابد که عدم افت کیفیت بالینی را تضمین می‌کند \lr{[4, 7]}.\par

% SOURCE Objectives_Chapter1.txt:23-23 [spacing]


% SOURCE Objectives_Chapter1.txt:24-24 [paragraph]
۵. ارزیابی کمّی مقاومت سیستم در برابر حملات تفاضلی و دستیابی به مقادیر \lr{NPCR} بالاتر از \lr{99.6\%} و \lr{UACI} حدود \lr{33.5\%} به همراه زمان اجرای زیر \lr{3} ثانیه برای کاربردهای زمان‌واقعی \lr{[1, 3]}.\par

% SOURCE Objectives_Chapter1.txt:25-25 [paragraph]
توضیح و بسط تفصیلی: این هدف فرعی بر سنجش‌پذیری و ارزیابی عملیاتی سامانه در برابر حملات تفاضلی و الزامات پردازش زمان‌واقعی \lr{(Real{-}time)} متمرکز است \lr{[1, 2, 3]}. شاخص نرخ تغییر تعداد پیکسل‌ها \lr{(NPCR)} با مقدار هدف بالای \lr{99.62\%} و شاخص میانگین شدت تغییرات مطلق \lr{(UACI)} با مقدار هدف \lr{33.51\%} نشان می‌دهند که کوچک‌ترین تغییر در تصویر ورودی منجر به تغییر کاملاً تصادفی در کل تصویر رمزشده می‌شود \lr{[1, 3]}. همچنین با بهینه‌سازی الگوریتم‌ها، زمان کل پردازش برای تصاویر با ابعاد \lr{256x256} به زیر ۳ ثانیه (حدود \lr{2.93} ثانیه) تقلیل می‌یابد که استفاده از این سامانه را در تجهیزات لبه \lr{IoMT} و بسترهای کم‌پهنای‌باند سلامت همراه میسر می‌سازد \lr{[1, 2]}.\par

% SOURCE Objectives_Chapter1.txt:26-26 [spacing]

% SOURCE Questions_Hypotheses_Chapter1.txt:1-1 [heading]
\SourceHeading{1}{۱{-}۴. پرسش‌ها و فرضیات پژوهش}

% SOURCE Questions_Hypotheses_Chapter1.txt:2-2 [spacing]


% SOURCE Questions_Hypotheses_Chapter1.txt:3-3 [paragraph]
برای دستیابی به هدف اصلی پژوهش و پاسخگویی به چالش‌های سه‌گانه مطرح‌شده در بیان مسئله، پرسش‌ها و فرضیات پژوهش به صورت متناظر در سه محور ساختاری تنظیم شده‌اند \lr{[1, 2, 7]}:\par

% SOURCE Questions_Hypotheses_Chapter1.txt:4-4 [spacing]


% SOURCE Questions_Hypotheses_Chapter1.txt:5-5 [heading]
\SourceHeading{2}{۱{-}۴{-}۱. پرسش‌های پژوهش}

% SOURCE Questions_Hypotheses_Chapter1.txt:6-6 [spacing]


% SOURCE Questions_Hypotheses_Chapter1.txt:7-7 [paragraph]
۱. پرسش اول (محور تولید کلید پویا و امنیت ساختاری): آیا تلفیق استخراج ویژگی‌های ناحیه \lr{ROI} توسط شبکه \lr{U{-}Net} با تابع هش، می‌تواند کلیدهای پویایی تولید کند که امنیت سیستم را در برابر حملات متن‌واضح برگزیده \lr{(CPA)} تضمین نماید؟ \lr{[1, 2, 9]}\par

% SOURCE Questions_Hypotheses_Chapter1.txt:8-8 [paragraph]
به بیان دقیق‌تر و تفصیلی، آیا استخراج اتصالات پرشی \lr{(Skip Connections)} و نقشه‌های ویژگی لایه‌های عمیق در معماری \lr{U{-}Net} به همراه محاسبه هش ۲۵۶ بیتی \lr{SHA{-}256} بر روی ناحیه حساس تشخیصی \lr{(ROI)}، قادر است کلید پویایی تولید نماید که با ایجاد اثر بهمنی \lr{(Avalanche Effect)} کامل، هرگونه وابستگی میان تصویر واضح و تصویر رمزنگاری‌شده را از بین برده و امکان نفوذ از طریق حملات متن‌واضح معلوم \lr{(KPA)} و متن‌واضح برگزیده \lr{(CPA)} را به طور کامل مسدود سازد؟ \lr{[1, 2]}.\par

% SOURCE Questions_Hypotheses_Chapter1.txt:9-9 [spacing]


% SOURCE Questions_Hypotheses_Chapter1.txt:10-10 [paragraph]
۲. پرسش دوم (محور آشفته‌سازی، انتشار و ارزیابی کمّی رمزنگاری): سیستم ابرآشوبی ۵ بعدی همراه با انتشار پویای \lr{DNA} تا چه حد قادر است همبستگی پیکسل‌های مجاور را شکسته و آنتروپی تصویر رمزنگاری‌شده را به حد آرمانی (۸ بیت) برساند؟ \lr{[1, 3, 6]}\par

% SOURCE Questions_Hypotheses_Chapter1.txt:11-11 [paragraph]
به عبارت دیگر، اعمال متوالی الگوریتم آشفته‌سازی زیگ‌زاگی \lr{(Zig{-}zag Scrambling)} برای اختلال موقعیت مکانی پیکسل‌ها و فازهای انتشار پویای \lr{DNA} (بر پایه عملیات پویای چرخش \lr{Dynamic DNA Flip} و \lr{XOR} آشوبی) تحت هدایت توالی‌های شبه‌تصادفی سیستم ۵ بعدی با فضای کلید فراتر از \lr{2}\SourceGlyph{005E}\lr{500}، تا چه اندازه می‌تواند مقادیر همبستگی پیکسلی در سه جهت افقی، عمودی و قطری را به صفر نزدیک کرده و معیارهای حساسیت تفاضلی \lr{NPCR} و \lr{UACI} را در محدوده استاندارد آرمانی (به ترتیب بالای \lr{99.62\%} و \lr{33.51\%}) قرار دهد؟ \lr{[1, 3, 6]}.\par

% SOURCE Questions_Hypotheses_Chapter1.txt:12-12 [spacing]


% SOURCE Questions_Hypotheses_Chapter1.txt:13-13 [paragraph]
۳. پرسش سوم (محور پنهان‌نگاری فراداده و حفظ کیفیت بالینی): آیا به‌کارگیری تکنیک پنهان‌نگاری مبتنی بر برجستگی بصری \lr{(Saliency)} امکان گنجاندن فراداده‌های بیمار بدون افت کیفیت تشخیصی \lr{(PSNR \SourceGlyph{003E} 50 dB)} را فراهم می‌سازد؟ \lr{[4, 5, 7]}\par

% SOURCE Questions_Hypotheses_Chapter1.txt:14-14 [paragraph]
به صورت تفکیکی، آیا الگوریتم پنهان‌نگاری آگاه به برجستگی بصری \lr{(Saliency{-}aware Steganography)} با اولویت‌دهی به نواحی کم‌اهمیت و پس‌زمینه تصویر و ترکیب آن با فشرده‌سازی اولیه \lr{zlib}، می‌تواند حجم بالایی از فراداده‌های هویتی و پرونده پزشکی بیمار را بدون ایجاد کوچک‌ترین اعوجاج دیداری در ناحیه حساس تشخیصی \lr{ROI} جای‌دهی نموده و شاخص ارزیابی \lr{PSNR} را بالای \lr{50} دسی‌بل و شاخص شباهت ساختاری \lr{SSIM} را در حد دست‌نخورده (نزدیک به \lr{1}) حفظ نماید؟ \lr{[4, 7, 8]}.\par

% SOURCE Questions_Hypotheses_Chapter1.txt:15-15 [spacing]


% SOURCE Questions_Hypotheses_Chapter1.txt:16-16 [spacing]


% SOURCE Questions_Hypotheses_Chapter1.txt:17-17 [heading]
\SourceHeading{2}{۱{-}۴{-}۲. فرضیات پژوهش}

% SOURCE Questions_Hypotheses_Chapter1.txt:18-18 [spacing]


% SOURCE Questions_Hypotheses_Chapter1.txt:19-19 [paragraph]
۱. فرضیه اول (مربوط به پرسش اول): به‌کارگیری شبکه \lr{U{-}Net} برای تولید کلید پویا از ناحیه \lr{ROI}، وابستگی کامل رمزنگاری به محتوای تصویر را ایجاد نموده و مقاومت سیستم در برابر حملات دیفرانسیلی و متن‌واضح معلوم را به طور قطعی افزایش می‌دهد \lr{[1, 2, 9]}.\par

% SOURCE Questions_Hypotheses_Chapter1.txt:20-20 [paragraph]
* تبیین تفصیلی فرضیه اول: استخراج خصوصیات آماری و پیکسل‌های کلیدی ناحیه \lr{ROI} از طریق بخش انقباضی و انبساطی \lr{U{-}Net} و اعمال تابع هش \lr{SHA{-}256} بر روی آن‌ها، کلید رمزنگاری را مستقیماً تابع محتوای همان تصویر می‌سازد. تغییر حتی یک پیکسل در تصویر ورودی منجر به تغییر بنیادی در حداقل \lr{50\%} از بیت‌های کلید نهایی گشته (اثر بهمنی) و در نتیجه شرایط اولیه سیستم ۵ بعدی کاملاً دگرگون می‌شود. این امر مقاومت الگوریتم را در برابر حملات متن‌واضح برگزیده \lr{(CPA)} و متن‌واضح معلوم \lr{(KPA)} تضمین می‌نماید \lr{[1, 2]}.\par

% SOURCE Questions_Hypotheses_Chapter1.txt:21-21 [paragraph]
  {-} متغیر مستقل: قطعه‌بندی \lr{ROI} با شبکه عمیق \lr{U{-}Net} و تولید کلید پویا بر پایه تابع هش \lr{SHA{-}256}.\par

% SOURCE Questions_Hypotheses_Chapter1.txt:22-22 [paragraph]
  {-} متغیر وابسته: میزان اثر بهمنی \lr{(Avalanche Effect)}، وابستگی رمز به تصویر ورودی و نرخ مقاومت در برابر حملات \lr{CPA} و \lr{KPA}.\par

% SOURCE Questions_Hypotheses_Chapter1.txt:23-23 [spacing]


% SOURCE Questions_Hypotheses_Chapter1.txt:24-24 [paragraph]
۲. فرضیه دوم (مربوط به پرسش دوم): الگوریتم آشفته‌سازی زیگ‌زاگی و انتشار پویای \lr{DNA} بر پایه سیستم ۵ بعدی، آنتروپی اطلاعات را به بالاتر از \lr{7.998} بیت رسانده و شاخص‌های \lr{NPCR} و \lr{UACI} را در محدوده استاندارد آرمانی (به ترتیب بالای \lr{99.6\%} و \lr{33.5\%}) قرار می‌دهد \lr{[1, 3, 6]}.\par

% SOURCE Questions_Hypotheses_Chapter1.txt:25-25 [paragraph]
* تبیین تفصیلی فرضیه دوم: ترکیب آشفته‌سازی زیگ‌زاگی (جهت نابودی همبستگی فضایی پیکسل‌ها) با فاز انتشار پویای \lr{DNA} بر پایه نگاشت‌های دینامیکی (چرخش و جمع پویای قواعد \lr{DNA})، تحت هدایت توالی‌های شبه‌تصادفی سیستم ابرآشوبی ۵ بعدی با توان‌های نمای لیاپانوف مثبت متعدد، هیستوگرام تصویر رمزنگاری‌شده را کاملاً یکنواخت می‌سازد. این فرآیند آنتروپی شانون را به \lr{7.9985} بیت (بسیار نزدیک به حد نظری \lr{8} بیت) رسانده و شاخص‌های حساسیت به تغییرات یک پیکسلی \lr{NPCR} را بالای \lr{99.62\%} و \lr{UACI} را در حدود \lr{33.51\%} ثبت می‌نماید \lr{[1, 3, 6]}.\par

% SOURCE Questions_Hypotheses_Chapter1.txt:26-26 [paragraph]
  {-} متغیر مستقل: توالی‌های شبه‌تصادفی سیستم ابرآشوبی ۵ بعدی، آشفته‌سازی زیگ‌زاگی و فازهای انتشار پویای \lr{DNA}.\par

% SOURCE Questions_Hypotheses_Chapter1.txt:27-27 [paragraph]
  {-} متغیر وابسته: ضریب همبستگی پیکسل‌های مجاور، آنتروپی اطلاعات شانون، و مقادیر آزمون‌های تفاضلی \lr{NPCR} و \lr{UACI}.\par

% SOURCE Questions_Hypotheses_Chapter1.txt:28-28 [spacing]


% SOURCE Questions_Hypotheses_Chapter1.txt:29-29 [paragraph]
۳. فرضیه سوم (مربوط به پرسش سوم): جاسازی فراداده‌های فشرده‌شده بیمار در نواحی کم‌برجستگی تصویر پزشکی، منجر به پنهان‌نگاری با کیفیت بالا \lr{(PSNR \SourceGlyph{003E} 50 dB)} شده و کیفیت بالینی و تشخیصی تصویر را تماماً حفظ می‌نماید \lr{[4, 7, 8]}.\par

% SOURCE Questions_Hypotheses_Chapter1.txt:30-30 [paragraph]
* تبیین تفصیلی فرضیه سوم: فشرده‌سازی فراداده‌های هویتی بیمار با الگوریتم \lr{zlib} باعث کاهش نرخ اشغال داده شده و تعیین نواحی کم‌برجستگی بصری (مبتنی بر نقشه \lr{Saliency Detection})، تضمین می‌کند که فرایند جای‌دهی داده‌ها تنها در نواحی فاقد ارزش تشخیصی (مانند پس‌زمینه سینه یا فضای اطراف مغز) صورت گیرد. این فرآیند منجر به دستیابی به نسبت سیگنال به نویز بیشینه \lr{PSNR} بالای \lr{53.4} دسی‌بل و شاخص شباهت ساختاری \lr{SSIM} برابر \lr{0.998} شده و کیفیت بالینی ناحیه حساس \lr{ROI} را بدون کوچک‌ترین اعوجاج حفظ می‌کند \lr{[4, 7, 8]}.\par

% SOURCE Questions_Hypotheses_Chapter1.txt:31-31 [paragraph]
  {-} متغیر مستقل: الگوریتم فشرده‌سازی \lr{zlib} و تکنیک پنهان‌نگاری آگاه به برجستگی بصری \lr{(Saliency Detection)}.\par

% SOURCE Questions_Hypotheses_Chapter1.txt:32-32 [paragraph]
  {-} متغیر وابسته: میزان نسبت سیگنال به نویز بیشینه \lr{(PSNR)}، شاخص شباهت ساختاری \lr{(SSIM)}، ظرفیت پنهان‌سازی و حفظ اصالت بالینی \lr{ROI}.\par

% SOURCE Questions_Hypotheses_Chapter1.txt:33-33 [spacing]

% SOURCE Innovations_Chapter1.txt:1-1 [heading]
\SourceHeading{1}{۱{-}۵. نوآوری‌های اصلی پژوهش}

% SOURCE Innovations_Chapter1.txt:2-2 [spacing]


% SOURCE Innovations_Chapter1.txt:3-3 [paragraph]
پژوهش حاضر دارای نوآوری‌های علمی، الگوریتمی و کاربردی متعددی در حوزه امنیت داده‌های سلامت دیجیتال است که با هدف غلبه بر چالش‌های همزمان رمزنگاری و پنهان‌نگاری تصاویر تشخیصی تدوین گردیده‌اند \lr{[1{-}10]}. مهم‌ترین جنبه‌های نوآوری و تمایز این پژوهش در چهار محور اساسی به شرح زیر تبیین می‌گردند:\par

% SOURCE Innovations_Chapter1.txt:4-4 [spacing]


% SOURCE Innovations_Chapter1.txt:5-5 [paragraph]
۱. تولید کلید رمزنگاری کاملاً پویا بر پایه هوش مصنوعی:\par

% SOURCE Innovations_Chapter1.txt:6-6 [paragraph]
تلفیق هوشمندانه معماری شبکه عصبی عمیق \lr{U{-}Net} و تابع هش \lr{SHA{-}256} بر روی ناحیه حساس تشخیصی \lr{(ROI)} تصویر پزشکی، راهکار نوآورانه‌ای برای تولید کلید پویای ۲۵۶ بیتی ارائه می‌دهد \lr{[1, 2]}. برخلاف سامانه‌های سنتی آشوبی که از کلیدهای ثابت، ایستا یا مستقل از تصویر ورودی استفاده می‌کنند، در این روش با کوچک‌ترین تغییر در تصویر ورودی (حتی تغییر ارزش یک پیکسل)، خروجی هش و در نتیجه شرایط اولیه و پارامترهای کنترل سیستم ابرآشوبی به صورت بنیادی تغییر می‌یابند. این امر اثر بهمنی \lr{(Avalanche Effect)} کاملی ایجاد نموده و مقاومت سیستم را در برابر حملات متن‌واضح برگزیده \lr{(CPA)} و متن‌واضح معلوم \lr{(KPA)} به طور معناداری ارتقا می‌دهد \lr{[1, 3, 5]}.\par

% SOURCE Innovations_Chapter1.txt:7-7 [spacing]


% SOURCE Innovations_Chapter1.txt:8-8 [paragraph]
۲. معرفی سیستم ابرآشوبی ۵ بعدی جدید با نگاشت‌های پویای \lr{DNA}:\par

% SOURCE Innovations_Chapter1.txt:9-9 [paragraph]
ارتقای بعد سیستم آشوبی به ۵ بعد و بهره‌گیری از بیش از یک توان نمای لیاپانوف مثبت \lr{(Multiple Positive Lyapunov Exponents)}، پیچیدگی غیرخطی فوق‌العاده‌ای را به همراه فضای کلید فراتر از ۲ به توان ۵۰۰ \lr{(2\SourceGlyph{005E}500)} ایجاد می‌نماید \lr{[1, 6]}. نوآوری اصلی این بخش، ترکیب توالی‌های شبه‌تصادفی سیستم ۵ بعدی با فازهای آشفته‌سازی زیگ‌زاگی \lr{(Zig{-}zag Scrambling)} و قوانین پویای چرخش و جمع \lr{DNA} \lr{(Dynamic DNA Flip and XOR)} است که برخلاف نگاشت‌های آشوبی کم‌بعد و ایستا، مقاومت بسیار بالایی در برابر بازسازی فضای فاز \lr{(Phase Space Reconstruction)} و حملات آماری و تفاضلی ایجاد می‌کند \lr{[1, 6, 8]}.\par

% SOURCE Innovations_Chapter1.txt:10-10 [spacing]


% SOURCE Innovations_Chapter1.txt:11-11 [paragraph]
۳. پنهان‌نگاری فراداده بیمار آگاه به برجستگی بصری \lr{(Saliency{-}aware Steganography)}:\par

% SOURCE Innovations_Chapter1.txt:12-12 [paragraph]
اولویت‌دهی هوشمند به نواحی کم‌توجه دیداری (مانند پس‌زمینه بی‌اهمیت در تصاویر \lr{MRI}، \lr{CT} و \lr{X{-}Ray}) بر پایه استخراج نقشه برجستگی بصری \lr{(Saliency Detection)}، نوآوری عمده‌ای برای جاسازی داده فراهم ساخته و کیفیت بالینی ناحیه حساس تشخیصی را دست‌نخورده باقی می‌گذارد \lr{[4, 7]}. در این ماژول، فراداده‌های هویتی و بالینی بیمار ابتدا با الگوریتم \lr{zlib} فشرده‌سازی شده و سپس صرفاً در نواحی کم‌اهمیت تصویر جای‌دهی می‌شوند؛ این تکنیک ظرفیت پنهان‌سازی بالا را همزمان با حفظ کیفیت دیداری تصویر (\lr{PSNR \SourceGlyph{003E} 50 dB} و \lr{SSIM} \SourceGlyph{007E} \lr{1}) محقق ساخته و ناحیه حساس تشخیصی \lr{ROI} را کاملاً بدون اعوجاج حفظ می‌نماید \lr{[4, 7, 10]}.\par

% SOURCE Innovations_Chapter1.txt:13-13 [spacing]


% SOURCE Innovations_Chapter1.txt:14-14 [paragraph]
۴. معماری یکپارچه و سبک با قابلیت پیاده‌سازی روی پردازنده‌های لبه \lr{(Edge Computing)} در سامانه‌های \lr{IoMT}:\par

% SOURCE Innovations_Chapter1.txt:15-15 [paragraph]
طراحی زنجیره پردازشی به گونه‌ای بهینه‌سازی شده است که ضمن حفظ بالاترین سطح امنیت و شکستن کامل همبستگی پیکسل‌ها در سه جهت افقی، عمودی و قطری، زمان اجرای کلی را به زیر ۳ ثانیه برای تصاویر \lr{256x256} کاهش دهد \lr{[2, 5]}. این ویژگی، معماری پیشنهادی را به یک راهکار عملیاتی، یکپارچه و سبک با قابلیت پیاده‌سازی بر روی پردازنده‌های کم‌مصرف لبه (مانند پلتفرم‌های سخت‌افزاری \lr{NVIDIA Jetson TX2}) در تجهیزات اینترنت اشیای پزشکی \lr{(IoMT)} و سامانه‌های سلامت همراه تبدیل می‌سازد \lr{[2, 5, 8]}.\par

% SOURCE Innovations_Chapter1.txt:16-16 [spacing]

% SOURCE Significance_Chapter1.txt:1-1 [heading]
\SourceHeading{1}{۱{-}۶. اهمیت و ضرورت پژوهش}

% SOURCE Significance_Chapter1.txt:2-2 [spacing]


% SOURCE Significance_Chapter1.txt:3-3 [paragraph]
با گسترش روزافزون فناوری‌های سلامت الکترونیک \lr{(E{-}Health)}، پزشکی از راه دور \lr{(Telemedicine)} و سامانه بایگانی و تبادل تصاویر پزشکی \lr{(PACS)}، تبادل الکترونیکی دادگان تشخیصی دیجیتال بر بستر شبکه‌های ناامن به یکی از ارکان لاینفک درمانی تبدیل گردیده است. با این حال، حفظ امنیت، یکپارچگی و حریم خصوصی این داده‌ها هنگام انتقال، نیازمند راهکارهای حفاظتی متمایز و استواری است. ضرورت انجام این پژوهش را می‌توان در چهار بعد اصلی و متکامل به شرح زیر تبیین نمود:\par

% SOURCE Significance_Chapter1.txt:4-4 [spacing]


% SOURCE Significance_Chapter1.txt:5-5 [paragraph]
۱. ضرورت اخلاقی و قانونی (حفظ حریم خصوصی بیمار): بر اساس قوانین بین‌المللی نظیر \lr{HIPAA} و \lr{GDPR}، صیانت از داده‌های هویتی و پزشکی بیماران ملزم به رعایت بالاترین استانداردهای امنیتی است. افشای تصاویر پزشکی می‌تواند تبعات حقوقی و اجتماعی سنگینی به همراه داشته باشد \lr{[2, 9]}. تصاویر پزشکی متداول اغلب در فرمت استاندارد \lr{DICOM} ذخیره‌سازی و منتقل می‌شوند که در ساختار خود حاوی فراداده‌های حساس \lr{EXIF} (شامل نام و نام خانوادگی بیمار، کد ملی، سن، تاریخچه بیماری و نظرات تشخیصی پزشک معالج) می‌باشند \lr{[7, 9]}. شنود، افشا یا دستکاری غیرمجاز این فراداده‌ها در مسیر انتقال، علاوه بر نقض صریح حریم خصوصی بیماران، خسارات اخلاقی و جرایم سنگین مالی و حقوقی را برای بسترهای درمانی به همراه دارد. از این رو، ترکیب الگوریتم‌های پنهان‌نگاری فراداده با رمزنگاری تصویر تشخیصی، یک ضرورت حتمی اخلاقی، حقوقی و قانونی در سامانه‌های سلامت هوشمند محسوب می‌گردد \lr{[2, 4, 7]}.\par

% SOURCE Significance_Chapter1.txt:6-6 [spacing]


% SOURCE Significance_Chapter1.txt:7-7 [paragraph]
۲. ضرورت فنی و علمی (محدودیت روش‌های کلاسیک): ناکارآمدی رمزهای سنتی متنی در برابر همبستگی بالای تصاویر پزشکی، ایجاب می‌کند که الگوریتم‌های نوینی مبتنی بر نظریه آشوب و یادگیری عمیق توسعه یابند \lr{[1, 3]}. تصاویر تشخیصی پزشکی دارای ویژگی‌های ذاتاً متمایزی نظیر حجم سنگین داده، افزونگی اطلاعاتی بسیار بالا و همبستگی شدید پیکسلی (با ضریب همبستگی نزدیک به ۱ در جهت‌های افقی، عمودی و قطری) می‌باشند \lr{[1, 3, 6]}. الگوریتم‌های متعارف رمزنگاری متنی نظیر استاندارد رمزنگاری پیشرفته \lr{(AES)} و \lr{RSA} به دلیل پیچیدگی محاسباتی بالا، ایجاد گلوگاه‌های پردازشی زمانی در تصاویر حجیم و عدم توانایی در شکستن همبستگی‌های پیکسلی دو‌بعدی، قادر به تامین امنیت تصاویر پزشکی در برنامه‌های کاربردی زمان‌واقعی نیستند \lr{[3, 6]}. همچنین سامانه‌های آشوبی کم‌بعد سنتی (مانند نگاشت‌های یک‌بعدی و دوبعدی) به دلیل فضای کلید محدود (کمتر از \lr{2}\SourceGlyph{005E}\lr{128}) و دوره تناوب کوتاه، در برابر حملات بازسازی فضای فاز و حملات دیفرانسیلی آسیب‌پذیرند \lr{[3]}. از این رو، الگوریتم‌های نوینی مبتنی بر یادگیری عمیق (معماری \lr{U{-}Net}) جهت تولید کلیدهای پویا و سیستم‌های ابرآشوبی ۵ بعدی با توان‌های لیاپانوف مثبت متعدد و انتشار \lr{DNA}، یک ضرورت علمی انکارناپذیر برای پاسخگویی به چالش‌های فنی رمزنگاری چندرسانه‌ای به شمار می‌روند \lr{[1, 2, 6]}.\par

% SOURCE Significance_Chapter1.txt:8-8 [spacing]


% SOURCE Significance_Chapter1.txt:9-9 [paragraph]
۳. ضرورت کاربردی در پزشکی از راه دور \lr{(Telemedicine)}: با توجه به توسعه درمان‌های از راه دور در مناطق محروم، ارسال امن و پرسرعت تصاویر تشخیصی بر روی شبکه‌های اینترنت همراه و کم‌پهنای‌باند نیازمند سامانه‌های رمزنگاری سریع و مقاوم است \lr{[2, 5]}. در بسترهای نوین سلامت همراه \lr{(m{-}Health)} و اینترنت اشیای پزشکی \lr{(IoMT)}، داده‌های پزشکی معمولاً توسط تجهیزات کم‌پهنای‌باند و دستگاه‌های لبه \lr{(Edge Devices)} جمع‌آوری و ارسال می‌شوند. اگر الگوریتم امنیتی سبک و زمان‌واقعی نباشد، تأخیر در انتقال تصاویر تشخیصی حاد (مانند \lr{MRI} مغزی یا \lr{CT} قفسه سینه) می‌تواند فرایند امدادرسانی بالینی را مختل سازد \lr{[2, 8]}. ضرورت کاربردی پژوهش حاضر در ارائه یک چارچوب امنیتی سبک و یکپارچه با زمان اجرای کلی زیر ۳ ثانیه برای تصاویر \lr{256x256} است که امکان پیاده‌سازی عملیاتی بر روی پردازنده‌های کم‌مصرف لبه (مانند \lr{NVIDIA Jetson TX2}) را فراهم آورده و امنیت تبادلات داده‌ای در بسترهای سلامت همراه را تضمین می‌کند \lr{[2, 5, 8]}.\par

% SOURCE Significance_Chapter1.txt:10-10 [spacing]


% SOURCE Significance_Chapter1.txt:11-11 [paragraph]
۴. ضرورت حفظ اصالت و کیفیت تشخیصی بالینی \lr{(Diagnostic Integrity)}: یکی از چالش‌های بنیادین در به‌کارگیری پنهان‌نگاری و رمزنگاری تصاویر پزشکی، عدم ایجاد نویز دیداری و اعوجاج در نواحی حساس آناتومیک \lr{(ROI)} است تا فرایند تشخیص دقیق توسط پزشک معالج دچار دگرگونی نگردد \lr{[4, 7]}. ضرورت این پژوهش در ارائه یک الگوریتم پنهان‌نگاری تطبیقی آگاه به برجستگی بصری \lr{(Saliency{-}aware Steganography)} نهفته است که با بهره‌گیری از فشرده‌سازی \lr{zlib}، فراداده‌های بیمار را صرفاً در نواحی کم‌اهمیت پس‌زمینه جای‌دهی می‌کند؛ این امر ضمن تامین ظرفیت پنهان‌سازی بالا، کیفیت دیداری تصویر را با \lr{PSNR} بالاتر از ۵۰ دسی‌بل و \lr{SSIM} نزدیک به ۱ حفظ کرده و اصالت بالینی ناحیه حساس تشخیصی را کاملاً دست‌نخورده باقی می‌گذارد \lr{[4, 7, 8]}.\par

% SOURCE Significance_Chapter1.txt:12-12 [spacing]

% SOURCE Scope_Chapter1.txt:1-1 [heading]
\SourceHeading{1}{۱{-}۷. حدود و قلمرو پژوهش}

% SOURCE Scope_Chapter1.txt:2-2 [spacing]


% SOURCE Scope_Chapter1.txt:3-3 [paragraph]
تعیین دقیق حدود و قلمرو پژوهش \lr{(Delimitations)}، مرزهای آگاهانه و انتخاب‌های روش‌شناختی پژوهشگر را جهت تمرکز بر اهداف اصلی، پرهیز از پراکنده‌گویی و تضمین تکرارپذیری آزمایش‌ها مشخص می‌سازد \lr{[1, 2]}. قلمرو این پژوهش در سه بخش موضوعی، داده‌ای و ابزاری به شرح زیر تفکیک و تبیین می‌گردد:\par

% SOURCE Scope_Chapter1.txt:4-4 [spacing]


% SOURCE Scope_Chapter1.txt:5-5 [paragraph]
۱. قلمرو موضوعی (محتوایی و روش‌شناختی):\par

% SOURCE Scope_Chapter1.txt:6-6 [paragraph]
تمرکز اصلی این پژوهش بر حوزه امنیت تصویر، رمزنگاری آشوبی، یادگیری عمیق (معماری \lr{U{-}Net}) و پنهان‌نگاری تصاویر پزشکی استوار است \lr{[1, 2, 7]}. به طور مشخص، مرزهای موضوعی شامل موارد زیر می‌باشد:\par

% SOURCE Scope_Chapter1.txt:7-7 [paragraph]
{-} قطعه‌بندی معنایی ناحیه حساس تشخیصی \lr{(ROI)} توسط شبکه عصبی عمیق \lr{U{-}Net} و استخراج ویژگی‌های آماری وابسته به محتوا جهت محاسبه کلید پویای ۲۵۶ بیتی با تابع هش \lr{SHA{-}256} \lr{[1, 2]}.\par

% SOURCE Scope_Chapter1.txt:8-8 [paragraph]
{-} فرمول‌بندی و شبیه‌سازی سیستم ابرآشوبی ۵ بعدی جدید با بیش از یک توان نمای لیاپانوف مثبت و اجرای فازهای متوالی آشفته‌سازی زیگ‌زاگی \lr{(Zig{-}zag Scrambling)} و انتشار پویای \lr{DNA} \lr{(Dynamic DNA Flip and XOR)} \lr{[1, 6]}.\par

% SOURCE Scope_Chapter1.txt:9-9 [paragraph]
{-} پنهان‌نگاری فراداده‌های هویتی و بالینی بیمار مبتنی بر استخراج نقشه برجستگی بصری \lr{(Saliency Detection)} و فشرده‌سازی اولیه با الگوریتم \lr{zlib} در تصاویر پزشکی ایستا (دو‌بعدی) \lr{[4, 7]}.\par

% SOURCE Scope_Chapter1.txt:10-10 [paragraph]
{-} استثنائات و مرزهای عدم شمول \lr{(Delimitations)}: پردازش ویدیوهای پزشکی زمان‌واقعی متحرک، دادگان ۳ بعدی پیوسته، پیاده‌سازی مستقیم در قالب تراشه‌های سخت‌افزاری اختصاصی \lr{(ASIC)} و الگوریتم‌های رمزنگاری کوانتومی خارج از حدود تمرکز پژوهش حاضر بوده و به عنوان زمینه‌های پیشنهادی برای پژوهش‌های آتی مطرح شده‌اند.\par

% SOURCE Scope_Chapter1.txt:11-11 [spacing]


% SOURCE Scope_Chapter1.txt:12-12 [paragraph]
۲. قلمرو داده‌ای (تنوع مدالیته‌ها و مجموعه‌داده‌های مرجع):\par

% SOURCE Scope_Chapter1.txt:13-13 [paragraph]
در این پژوهش از ۴ مجموعه‌داده مرجع بین‌المللی شامل تصاویر شبکیه چشم (\lr{DRIVE} و \lr{RITE})، تصاویر \lr{MRI} مغزی \lr{(BraTS)} و تصاویر رادیوگرافی قفسه سینه \lr{(COVID{-}19 Chest X{-}Ray)} استفاده شده است \lr{[1, 2, 5]}. توجیه روش‌شناختی انتخاب این دادگان عبارت است از:\par

% SOURCE Scope_Chapter1.txt:14-14 [paragraph]
{-} اثبات تعمیم‌پذیری و مقاومت چارچوب پیشنهادی بر روی سه مدالیته متداول و حیاتی تصویربرداری پزشکی (چشم‌پزشکی، نورولوژی و بیماری‌های ریوی) با ساختارهای بافتی، ابعاد، و همبستگی‌های پیکسلی متفاوت.\par

% SOURCE Scope_Chapter1.txt:15-15 [paragraph]
{-} مجموعه‌داده‌های شبکیه چشم (\lr{DRIVE} و \lr{RITE}): جهت ارزیابی کارایی \lr{U{-}Net} در قطعه‌بندی دقیق عروق خونی باریک و پیچیده و سنجش رمزنگاری پاتولوژی‌های دقیق چشمی \lr{[1, 2]}.\par

% SOURCE Scope_Chapter1.txt:16-16 [paragraph]
{-} مجموعه‌داده \lr{MRI} مغزی \lr{(BraTS)}: جهت سنجش عملکرد الگوریتم در بافت‌های نرم مغزی، قطعه‌بندی تومورها و ارزیابی ماژول پنهان‌نگاری فراداده در نواحی پس‌زمینه کم‌توجه \lr{[7]}.\par

% SOURCE Scope_Chapter1.txt:17-17 [paragraph]
{-} مجموعه‌داده رادیوگرافی قفسه سینه \lr{(Chest X{-}Ray)}: جهت سنجش پایداری و آنتروپی اطلاعات بر روی دادگان با افزونگی اطلاعاتی و همبستگی پیکسلی بسیار بالا \lr{[5, 8]}.\par

% SOURCE Scope_Chapter1.txt:18-18 [spacing]


% SOURCE Scope_Chapter1.txt:19-19 [paragraph]
۳. قلمرو ابزاری و پیاده‌سازی:\par

% SOURCE Scope_Chapter1.txt:20-20 [paragraph]
شبیه‌سازی و ارزیابی کمّی الگوریتم‌ها در محیط نرم‌افزاری \lr{Python 3.12} (با کتابخانه‌های \lr{PyTorch} و \lr{OpenCV}) و \lr{MATLAB R2024b} اجرا گردیده است \lr{[1, 3]}. تفکیک وظایف این ابزارها به شرح زیر است:\par

% SOURCE Scope_Chapter1.txt:21-21 [paragraph]
{-} محیط نرم‌افزاری \lr{Python 3.12}: به‌کارگیری کتابخانه تخصصی \lr{PyTorch} جهت پیاده‌سازی و آموزش شبکه \lr{U{-}Net} در استخراج \lr{ROI}، کتابخانه \lr{OpenCV} و \lr{NumPy} جهت پردازش تصویر و اجرای فازهای آشفته‌سازی زیگ‌زاگی و انتشار پویای \lr{DNA}، و ماژول \lr{zlib} جهت فشرده‌سازی فراداده‌ها.\par

% SOURCE Scope_Chapter1.txt:22-22 [paragraph]
{-} محیط نرم‌افزاری \lr{MATLAB R2024b}: جهت تحلیل رفتار غیرخطی و دینامیکی سیستم ۵ بعدی، محاسبه توان‌های نمای لیاپانوف، رسم پرتره‌های فاز ۵ بعدی (جذابه‌های غریب) و اجرای آزمون‌های آماری سنجش تصادفی‌بودن \lr{(NIST)} \lr{[1, 3, 6]}.\par

% SOURCE Scope_Chapter1.txt:23-23 [paragraph]
{-} ارزیابی محاسباتی و سخت‌افزاری: شبیه‌سازی بر روی پردازنده گرافیکی \lr{(GPU)} انجام شده و زمان اجرای کلی سیستم (زیر ۳ ثانیه) با الزامات سخت‌افزارهای لبه کم‌مصرف (مانند \lr{NVIDIA Jetson TX2}) در تجهیزات اینترنت اشیای پزشکی \lr{(IoMT)} سنجیده و تطبیق داده شده است \lr{[2, 8]}.\par

% SOURCE Scope_Chapter1.txt:24-24 [spacing]

% SOURCE Definitions_Chapter1.txt:1-1 [heading]
\SourceHeading{1}{۱{-}۸. تعاریف مفاهیم و واژگان تخصصی}

% SOURCE Definitions_Chapter1.txt:2-2 [spacing]


% SOURCE Definitions_Chapter1.txt:3-3 [paragraph]
برای ابهام‌زدایی، تسهیل در درک مفاهیم کلیدی، ایجاد زبان علمی مشترک و تبیین دقیق متغیرهای پژوهش، تعاریف مفاهیم و واژگان تخصصی در دو سطح تفکیک‌شده ارائه می‌گردند: تعاریف نظری (مفهوم‌پردازی پایه بر اساس ادبیات پژوهشی و مراجع مرجع) و تعاریف عملیاتی (نحوه پیاده‌سازی، فرمول‌بندی و سنجش عددی متغیرها در این پایان‌نامه) \lr{[1{-}10]}.\par

% SOURCE Definitions_Chapter1.txt:4-4 [spacing]


% SOURCE Definitions_Chapter1.txt:5-5 [heading]
\SourceHeading{2}{۱{-}۸{-}۱. تعاریف نظری}

% SOURCE Definitions_Chapter1.txt:6-6 [spacing]


% SOURCE Definitions_Chapter1.txt:7-7 [paragraph]
۱. سیستم ابرآشوبی \lr{(Hyperchaotic System)}: یک سیستم دینامیکی غیرخطی با حداقل چهار بعد در فضای فاز است که دارای بیش از یک توان نمای لیاپانوف مثبت \lr{(Multiple Positive Lyapunov Exponents)} می‌باشد. این سیستم‌ها رفتاری بسیار پیچیده‌تر، تصادفی‌تر و غیرقابل‌پیش‌بینی‌تر از سیستم‌های آشوبی معمولی (مانند نگاشت‌های یک‌بعدی یا دوبعدی) نشان داده و پایداری فوق‌العاده‌ای در برابر حملات بازسازی فضای فاز \lr{(Phase Space Reconstruction)} ارائه می‌دهند \lr{[1, 3, 6]}.\par

% SOURCE Definitions_Chapter1.txt:8-8 [spacing]


% SOURCE Definitions_Chapter1.txt:9-9 [paragraph]
۲. ناحیه‌بندی معنایی \lr{(Semantic Segmentation)}: فرآیند دسته‌بندی و برچسب‌زنی دقیق هر پیکسل از تصویر به یک کلاس خاص آناتومیک یا پاتولوژیک (مانند ضایعه، تومور، عروق خونی یا بافت سالم) در شبکه عصبی عمیق است، به گونه‌ای که مرزهای ساختاری نواحی حساس تشخیصی با دقت بالا از زمینه تفکیک گردند \lr{[1, 2]}.\par

% SOURCE Definitions_Chapter1.txt:10-10 [spacing]


% SOURCE Definitions_Chapter1.txt:11-11 [paragraph]
۳. پنهان‌نگاری \lr{(Steganography)}: هنر و دانش پنهان‌سازی اطلاعات محرمانه (مانند فراداده هویتی و پرونده بالینی بیمار) در درون یک رسانه پوششی (مانند تصویر پزشکی) به گونه‌ای که وجود پیام مخفی از دید نفوذگران و ابزارهای پنهان‌کاوی \lr{(Steganalysis)} کشف‌نشده و غیرقابل‌تشخیص باقی بماند \lr{[4, 7, 8]}.\par

% SOURCE Definitions_Chapter1.txt:12-12 [spacing]


% SOURCE Definitions_Chapter1.txt:13-13 [paragraph]
۴. برجستگی بصری \lr{(Saliency Detection)}: مکانیزمی مبتنی بر الگوریتم‌های بینایی ماشین و هوش مصنوعی که نواحی پرتوجه و دارای ارزش اطلاعاتی دیداری بالایی (مانند بافت‌های تشخیصی) را از نواحی کم‌اهمیت و پس‌زمینه تصویر تفکیک می‌نماید \lr{[5, 7]}.\par

% SOURCE Definitions_Chapter1.txt:14-14 [spacing]


% SOURCE Definitions_Chapter1.txt:15-15 [paragraph]
۵. تابع هش یک‌طرفه \lr{(One{-}Way Hash Function)} و اثر بهمنی \lr{(Avalanche Effect)}: تابع ریاضی یک‌طرفه‌ای که داده‌هایی با طول متغیر را به یک رشته بیتی ثابت (مانند ۲۵۶ بیت در \lr{SHA{-}256}) تبدیل می‌کند؛ به طوری که محاسبه معکوس آن غیرممکن بوده و با کوچک‌ترین تغییر در تصویر ورودی (حتی تغییر یک پیکسل)، حداقل ۵۰\SourceGlyph{066A} از بیت‌های خروجی هش تغییر می‌یابند که به این خصوصیت اثر بهمنی گفته می‌شود \lr{[1, 2]}.\par

% SOURCE Definitions_Chapter1.txt:16-16 [spacing]


% SOURCE Definitions_Chapter1.txt:17-17 [paragraph]
۶. فشرده‌سازی اطلاعات \lr{(zlib Data Compression)}: الگوریتم فشرده‌سازی داده‌های دیجیتال بدون افت \lr{(Lossless)} که بر پایه تلفیق الگوریتم‌های \lr{LZ77} و کدگذاری هفمن عمل کرده و حجم فراداده‌های متنی را قبل از فاز پنهان‌نگاری کاهش می‌دهد \lr{[7, 8]}.\par

% SOURCE Definitions_Chapter1.txt:18-18 [spacing]


% SOURCE Definitions_Chapter1.txt:19-19 [paragraph]
۷. شاخص شباهت ساختاری \lr{(SSIM)}: معیاری دیداری و آماری جهت سنجش میزان شباهت ساختاری، روشنایی و کنتراست میان دو تصویر (تصویر اصلی و تصویر پنهان‌نگاری‌شده) که دامنه تغییرات آن بین ۰ تا ۱ قرار دارد \lr{[4, 7, 8]}.\par

% SOURCE Definitions_Chapter1.txt:20-20 [spacing]


% SOURCE Definitions_Chapter1.txt:21-21 [heading]
\SourceHeading{2}{۱{-}۸{-}۲. تعاریف عملیاتی}

% SOURCE Definitions_Chapter1.txt:22-22 [spacing]


% SOURCE Definitions_Chapter1.txt:23-23 [paragraph]
۱. شبکه \lr{U{-}Net}: یک معماری شبکه عصبی کانولوشنی متقارن شامل مسیرهای انقباضی \lr{(Encoder)} و انبساطی \lr{(Decoder)} با اتصالات پرشی \lr{(Skip Connections)}. در این پژوهش، شبکه \lr{U{-}Net} جهت قطعه‌بندی خودکار و استخراج دقیق ناحیه حساس تشخیصی \lr{(ROI)} از دادگان مرجع \lr{DRIVE}، \lr{RITE}، \lr{BraTS} و \lr{Chest X{-}Ray} آموزش دیده و ویژگی‌های آماری متمایز این ناحیه جهت محاسبه کلید پویا به کار می‌رود \lr{[1, 2]}.\par

% SOURCE Definitions_Chapter1.txt:24-24 [spacing]


% SOURCE Definitions_Chapter1.txt:25-25 [paragraph]
۲. کلید پویای تصویر{-}محور \lr{(Image{-}Dependent Dynamic Key)}: در این پژوهش، کلید رمزنگاری ۲۵۶ بیتی حاصل از اعمال تابع هش \lr{SHA{-}256} بر روی خصوصیات آماری و ماتریس \lr{ROI} خروجی شبکه \lr{U{-}Net} است. این کلید به‌صورت پویا، شرایط اولیه \lr{(Variables)} و پارامترهای کنترل سیستم ابرآشوبی ۵ بعدی را مقداردهی اولیه کرده و وابستگی ۱۰۰ درصدی رمزنگاری به محتوای تصویر ورودی را تضمین می‌نماید \lr{[1, 2, 5]}.\par

% SOURCE Definitions_Chapter1.txt:26-26 [spacing]


% SOURCE Definitions_Chapter1.txt:27-27 [paragraph]
۳. آنتروپی شانون \lr{(Shannon Entropy)}: معیاری برای سنجش تصادفی بودن، عدم قطعیتی اطلاعات و توزیع یکنواخت سطوح خاکستری در تصویر رمزنگاری‌شده که مقدار آرمانی آن برای تصویر ۸ بیتی برابر ۸ می‌باشد. در این پژوهش، آنتروپی بر روی تصاویر رمزشده محاسبه گردیده و دست‌یابی به مقدار متوسط \lr{7.9985} بیت به عنوان نشانه تصادفی بودن کامل خروجی ارزیابی می‌گردد \lr{[1, 3, 6]}.\par

% SOURCE Definitions_Chapter1.txt:28-28 [spacing]


% SOURCE Definitions_Chapter1.txt:29-29 [paragraph]
۴. شاخص‌های \lr{NPCR} و \lr{UACI}: معیارهای ارزیابی کمّی میزان حساسیت الگوریتم رمزنگاری در برابر تغییرات یک پیکسلی در تصویر واضح (آزمون‌های مقاومت در برابر حملات تفاضلی). در این پژوهش، مقادیر عددی هدف برای \lr{NPCR} (تعداد پیکسل‌های تغییریافته) بالاتر از \lr{99.62\%} و برای \lr{UACI} (میانگین شدت تغییرات) حدود \lr{33.51\%} تنظیم و شبیه‌سازی شده‌اند که مطابق با حدود آرمانی تئوریک می‌باشند \lr{[1, 3]}.\par

% SOURCE Definitions_Chapter1.txt:30-30 [spacing]


% SOURCE Definitions_Chapter1.txt:31-31 [paragraph]
۵. نسبت سیگنال به نویز بیشینه \lr{(PSNR)} و \lr{SSIM}: معیارهایی بر حسب دسی‌بل و ضریب ساختاری جهت سنجش کیفیت دیداری تصویر پنهان‌نگاری‌شده نسبت به تصویر اصلی. در این پژوهش، نیل به \lr{PSNR} بالاتر از \lr{50} دسی‌بل (مقدار تجربی \lr{53.4 dB}) و \lr{SSIM} برابر \lr{0.998} به عنوان معیار عملیاتی حفظ کامل کیفیت بالینی و عدم ایجاد اعوجاج در ناحیه تشخیصی \lr{ROI} ارزیابی می‌شود \lr{[4, 7, 8]}.\par

% SOURCE Definitions_Chapter1.txt:32-32 [spacing]


% SOURCE Definitions_Chapter1.txt:33-33 [paragraph]
۶. پنهان‌نگاری آگاه به برجستگی \lr{(Saliency{-}aware Steganography)}: در این پژوهش، فراداده‌های هویتی و بالینی بیمار پس از فشرده‌سازی با الگوریتم \lr{zlib}، صرفاً در پیکسل‌های متعلق به نواحی کم‌برجستگی بصری (پس‌زمینه بی‌اهمیت تشخیصی) جای‌دهی می‌شوند؛ به گونه‌ای که ظرفیت پنهان‌سازی مناسب همراه با دست‌نخورده ماندن کامل کیفیت بالینی ناحیه \lr{ROI} محقق گردد \lr{[4, 7]}.\par

% SOURCE Definitions_Chapter1.txt:34-34 [spacing]

% SOURCE Structure_Chapter1.txt:1-1 [heading]
\SourceHeading{1}{۱{-}۹. ساختار عمومی پایان‌نامه}

% SOURCE Structure_Chapter1.txt:2-2 [spacing]


% SOURCE Structure_Chapter1.txt:3-3 [paragraph]
این پایان‌نامه جهت دست‌یابی به اهداف پژوهش و پاسخگویی به فرضیات و پرسش‌های مطرح‌شده، در قالب پنج فصل منطقی، منسجم و مرتبط به شرح زیر سازمان‌دهی شده است \lr{[\rl{\persianfont ۱}{-}\rl{\persianfont ۱۰}]}:\par

% SOURCE Structure_Chapter1.txt:4-4 [spacing]


% SOURCE Structure_Chapter1.txt:5-5 [paragraph]
فصل اول (کلیات پژوهش): شامل تبیین مقدماتی، بیان مسئله و چالش‌های امنیتی موجود در انتقال تصاویر پزشکی، اهداف اصلی و فرعی کمّی، پرسش‌ها و فرضیات پژوهش، نوآوری‌های چهارگانه علمی و الگوریتمی، اهمیت و ضرورت پژوهش در ابعاد اخلاقی، فنی و کاربردی، حدود و قلمرو داده‌ای و ابزاری، و در نهایت تعاریف نظری و عملیاتی متغیرها می‌باشد. این فصل به عنوان نقشه راه و شالوده مفهومی پایان‌نامه، زیربنای ورود به بررسی‌های تخصصی را فراهم می‌سازد.\par

% SOURCE Structure_Chapter1.txt:6-6 [spacing]


% SOURCE Structure_Chapter1.txt:7-7 [paragraph]
فصل دوم (مبانی نظری و پیشینه پژوهش): مرور ادبیات تحقیق، مفاهیم ابرآشوب، یادگیری عمیق، الگوریتم‌های \lr{DNA}، پنهان‌نگاری و بررسی نقادانه کارهای پیشین. در بخش نخست این فصل، اصول ریاضی سیستم‌های دینامیکی غیرخطی و توان‌های نمای لیاپانوف، معماری شبکه‌های عصبی عمیق کانولوشنی (به‌ویژه \lr{U{-}Net})، پویایی‌های کدگذاری \lr{DNA}، و تکنیک‌های بینایی ماشین در تشخیص برجستگی بصری \lr{(Saliency Detection)} تشریح می‌گردند. در بخش دوم، پژوهش‌های گذشته در حوزه رمزنگاری و پنهان‌نگاری تصاویر تشخیصی مورد نقد و ارزیابی قرار گرفته و شکاف‌های پژوهشی موجود استخراج می‌شوند \lr{[\rl{\persianfont ۱}{-}\rl{\persianfont ۱۰}]}.\par

% SOURCE Structure_Chapter1.txt:8-8 [spacing]


% SOURCE Structure_Chapter1.txt:9-9 [paragraph]
فصل سوم (روش پیشنهادی و معماری سیستم): تشریح دقیق معماری چهارلایه سیستم، استخراج کلید با \lr{U{-}Net}، فرمول‌بندی ۵ بعدی، فازهای جایگشت زیگ‌زاگی و انتشار \lr{DNA} و ماژول پنهان‌نگاری \lr{Saliency}. در این فصل، جزئیات الگوریتمی و ریاضی سیستم پیشنهادی شامل: (۱) ماژول قطعه‌بندی معنایی \lr{ROI} توسط \lr{U{-}Net} و استخراج کلید پویای ۲۵۶ بیتی \lr{SHA{-}256}، (۲) فرمول‌بندی معادلات دیفرانسیل سیستم ابرآشوبی ۵ بعدی جدید و محاسبه جذابه‌های غریب آن، (۳) فازهای آشفته‌سازی زیگ‌زاگی \lr{(Zig{-}zag Scrambling)} و انتشار پویای \lr{DNA} (شامل قوانین چرخش و جمع پویای \lr{XOR})، و (۴) ماژول پنهان‌نگاری فراداده بیمار مبتنی بر نقشه برجستگی بصری \lr{(Saliency{-}aware)} و فشرده‌سازی \lr{zlib} به تفصیل تشریح می‌گردند.\par

% SOURCE Structure_Chapter1.txt:10-10 [spacing]


% SOURCE Structure_Chapter1.txt:11-11 [paragraph]
فصل چهارم (طرح آزمایش‌ها و ارزیابی نتایج): معرفی دادگان، سناریوهای آزمایشی، تحلیل کمّی آنتروپی، همبستگی، \lr{NPCR}، \lr{UACI}، \lr{PSNR} و سنجش مقاومت در برابر حملات سایبری. این فصل به شبیه‌سازی و پیاده‌سازی عملیاتی طرح بر روی ۴ مجموعه‌داده مرجع \lr{(DRIVE, RITE, BraTS, COVID{-}19 Chest X{-}Ray)} اختصاص دارد. عملکرد سیستم در دو محیط پایتون و متلب بر پایه شاخص‌های آنتروپی شانون، ضریب همبستگی پیکسلی ۳جهته، آزمون‌های تفاضلی \lr{NPCR} و \lr{UACI}، آزمون‌های آماری \lr{NIST}، معیارهای کیفیت دیداری \lr{PSNR} و \lr{SSIM}، ارزیابی مقاومت در برابر حملات نویز و قیچی‌شدن، و تحلیل سرعت اجرا بر روی پردازنده‌های لبه \lr{IoMT} ارزیابی و با روش‌های مرجع مقایسه می‌گردد \lr{[\rl{\persianfont ۱}{-}\rl{\persianfont ۱۰}]}.\par

% SOURCE Structure_Chapter1.txt:12-12 [spacing]


% SOURCE Structure_Chapter1.txt:13-13 [paragraph]
فصل پنجم (جمع‌بندی، نتیجه‌گیری و پیشنهادها): پاسخ به پرسش‌ها و فرضیات، تبیین محدودیت‌ها و ارائه پیشنهادهای کاربردی برای پژوهش‌های آتی. در فصل پایانی، یافته‌های تجربی و کمّی پژوهش جمع‌بندی شده و پاسخ‌های صریح به پرسش‌ها و فرضیات سه‌گانه ارائه می‌شود. همچنین محدودیت‌های عملیاتی و اجرایی طرح تبیین گردیده و نقشه راه مشخصی در قالب پیشنهادهای علمی و فنی برای توسعه پژوهش‌های آتی در حوزه امنیت سلامت الکترونیک ترسیم می‌گردد.\par

% SOURCE Structure_Chapter1.txt:14-14 [spacing]

